\documentclass{article}
\usepackage{csquotes}
\usepackage{/globals/de}
\begin{document}
\section{Definition Verwaltungshandeln} 
Das Verwaltungshandeln beschreibt das Handeln der Verwaltung innerhalb eines Staates. Das Ziel der Verwaltung ist es dabei die ihr durch Gesetze aufgetragenen öffentlichen Aufgaben zu erfüllen.

Dies passiert durch das Treffen von Entscheidungen und dem Nehmen von Tätigkeiten und Maßnahmen durch die Behörden der Verwaltung.

Dazu gehören der Vollzug von Gesetzen, das Treffen von Ermessensentscheidungen, das Ausführen von Verwaltungsakten und Realakten, \usw.

\subsection{Polizei}
Im Sinne der Verwaltung und somit auch im Sinne dieses Moduls ist mit der \textquote{Polizei} nicht einfach die Institution der Polizei gemeint, sondern eher um dessen Handeln in Bezug auf dessen Aufgaben als Teil der Exekutive innerhalb des öffentlichen Rechts. Dazu gehört die Gefahrenabwehr und die Strafverfolgung
\begin{description}
 \item[Gefahrenabwehr] Das Verhindern oder Beseitigen von Gefahren, bereits bevor es zum Schaden kommt. 
 \item[Strafverfolgung] Die Aufklärung und Ahndung von Straftaten durch die Polizei, die Staatsanwaltschaft und Gerichte.
\end{description}
\end{document}